\t{Example 1}

Consider the following call:

\t{count\_tastiness(3, [5, 2, 1])}

This means that Aunty wants to pack $3$ bags, and there are $3$ types of biscuits in the shop:
\begin{itemize}
\item $5$ biscuits of type $0$, each having a tastiness value $1$,
\item $2$ biscuits of type $1$, each having a tastiness value $2$,
\item $1$ biscuit of type $2$, each having a tastiness value $4$.
\end{itemize}

The possible values of $y$ are $[0, 1, 2, 3, 4]$. For instance, in order to pack $3$ bags of total tastiness $3$, Aunty can pack:
\begin{itemize}
 \item one bag containing three biscuits of type $0$, and
 \item two bags, each containing one biscuit of type $0$ and one biscuit of type $1$. 
\end{itemize}

Since there are $5$ possible values of $y$, the procedure should return $5$.

\includegraphics[scale=0.8]{1.png}

\textbf{Example 2}

Consider the following call:

\t{count\_tastiness(2, [2, 1, 2])}

This means that Aunty wants to pack $2$ bags, and there are $3$ types of biscuits in the shop:
\begin{itemize}
\item $2$ biscuits of type $0$, each having a tastiness value $1$,
\item $1$ biscuits of type $1$, each having a tastiness value $2$,
\item $2$ biscuits of type $2$, each having a tastiness value $4$.
\end{itemize}

The possible values of $y$ are $[0,1,2,4,5,6]$. Since there are $6$ possible values of $y$, the procedure should return $6$.

